% !TeX root = ../../../main.tex
밤하늘의 별을 구경하던 기룡이는 직접 별자리를 만들려고 한다.
하지만 아무 별자리나 만드는 건 재미없기 때문에, 특별한 규칙을 적용하려고 한다.

\begin{itemize}
    \item 모든 별은 2차원 평면 위의 좌표 $(x, y)$와 1 이상 10 이하의 밝기 $b$로 표현되며, $b$가 8 이상인 별은 중심별이 된다.
    \item 별자리는 하나의 중심별과 그 중심별에 직접 또는 간접적으로 연결된 별들의 집합이다. 중심별 하나만으로도 별자리를 이룰 수 있다.
    \item 모든 별은 두 개 이상의 중심별과 직접 또는 간접적으로 연결될 수 없으며, 오직 하나의 중심별과만 연결되어야 한다. 따라서 중심별끼리는 연결할 수 없다.
    \item 두 별 $s$와 $t$를 연결하는 데 드는 비용은 $(b_{s}+b_{t})\times \left( (x_{s}-x_{t})^2+(y_{s}-y_{t})^2\right)$이다.
\end{itemize}

모든 별이 별자리에 포함되도록 연결할 때, 연결 비용의 합의 최솟값을 구해 보자.

좌표 $(x, y)$가 같은 별이 여러 개 주어지면, 밝기 $b$가 가장 큰 별 하나만 존재하는 것으로 간주한다.
